\documentclass[10pt]{article}

% Specify the margins and text size.
\setlength{\textwidth}{6.4in}
\setlength{\textheight}{9.5in}
\setlength{\oddsidemargin}{0pt}
\setlength{\evensidemargin}{0pt}
\setlength{\topmargin}{0pt}
\setlength{\hoffset}{.05in}
\setlength{\voffset}{-1in}

\setlength{\parskip}{5pt}
\setlength{\parindent}{0pt}

% Load some fonts and symbol packages
\usepackage{latexsym}
\usepackage{pifont}       % contains 'star' symbol for counterinsurgency handbook title
\usepackage{yfonts} 
\usepackage{amsmath}
\usepackage{amsfonts}

\usepackage{graphicx}     % actually, this is loaded by pstricks
\usepackage[T1]{fontenc}
\usepackage{ifthen}
\usepackage{pstricks,pst-grad,pst-text,pst-node,multido,pst-plot,calc,pst-3dplot}
%\usepackage[all]{xy}
%\usepackage{animate}

% The hyperref package inserts links into the pdf.
\definecolor{MyLinkColor}{rgb}{.1,.2,1}
\definecolor{MyCiteColor}{rgb}{.1,1,.2}
\definecolor{MyURLColor}{rgb}{.4,.4,.4}
\usepackage[backref=true,pagebackref=false,hyperindex,colorlinks=true,
  linkcolor=MyLinkColor,urlcolor=MyURLColor]{hyperref}


% The tweaklist package is something I found on the web.  It provides a simple interface
% for making changes to spacing used in the itemize and enumerate environments.  Comment
% this out if you don't care to use tweaklist.
\usepackage{tweaklist}
\renewcommand{\itemhook}{\setlength{\parskip}{2pt}\setlength{\parsep}%
{1pt}\setlength{\topsep}{0pt}\setlength{\itemsep}{0pt}}

\newcommand{\U}{\underline{\hspace{5pt}}}

\usepackage{listings}

\begin{document}
\pagestyle{empty}
\lstset{language=R, showspaces=false, showstringspaces=false}

\href{http://www.shepherd.edu}{\includegraphics[height=1.75cm]{logo-high-res.eps}}
\vspace{-1.69cm}

{\small
\begin{tabular}{cl}
& Math 314\\
& Statistics\\
\hspace{5.28in} &% 3 February 2012
\end{tabular}
}
\setlength{\baselineskip}{1.05\baselineskip}

\begin{center}
\textbf{\large  Correlation Example}
\end{center}
\newcommand{\Z}{\hphantom{0}}
\newcommand{\VP}{\vphantom{\Large Y}}
\newcommand{\M}{\hphantom{$-$}}

\begin{center}
\begin{tabular}{|rr|ccc||cc|}\hline
    &     &       &       &           & $y$ values   & \VP\\
    &     &       &       &           &  predicted  & SD\\
$x$ & $y$ & $z_x$ & $z_y$ & $z_x\,z_y$ & by SD line  & error\\[2pt]\hline
$1$   & $5$   & $-3/2$\M  & $-1/2$\M  & $3/4$       &  1  & 4\VP\\
$3$   & $9$   & $-1/2$\M  & $1/2$   & $-1/4$\M      &  5  & 4\\
$4$   & $7$   & $0$     & $0$     & $0$         &  7  & 0\\
$5$   & $1$   & $1/2$   & $-3/2$\M  & $-3/4$\M      &  9  & $-8$\M\\
$7$   & $13$  & $3/2$   & $3/2$   & $9/4$       &  13 & 0\\ \hline
\end{tabular}
\end{center}

\begin{itemize}
\item  Compute the necessary statistics.
\begin{align*}
\mbox{mean}(x) & = \frac{1}{5}\,(1 + 3 + 4 + 5 + 7) = 4\\
\mbox{SD}(x)   & = \sqrt{\frac{(1-4)^2 + (3-4)^2 + (4-4)^2 + (5-4)^2 + (7-4)^2}{5}} = 
   \sqrt{\frac{9 + 1 + 0 + 1 + 9}{5}} = \sqrt{\frac{20}{5}}=\sqrt{4} \\
               & =2\\
\mbox{mean}(y) &=\frac{1}{5}\,(5 + 9 + 7 + 1 + 13) = 7\\
\mbox{SD}(y)   &=\sqrt{\frac{(5-7)^2 + (9-7)^2 + (7-7)^2 + (1-7)^2 + (13-7)^2}{5}}
     = \sqrt{\frac{4 + 4 + 0 + 36 + 36}{5}} = \sqrt{\frac{80}{5}}\\
    & = 4
\end{align*}
\item Convert the $x$ values to standard units.  For example,
\[x = 1\mbox{\hspace{10pt}becomes\hspace{10pt}} z_x\frac{1-4}{2} = -3/2
\mbox{\hspace{20pt}and\hspace{20pt}}
x = 3\mbox{\hspace{10pt}becomes\hspace{10pt}} z_x\frac{3-4}{2} = -1/2\]
\item Convert the $y$ values to standard units:
For example,
\[y=5\mbox{\hspace{10pt}becomes\hspace{10pt}} z_y=\frac{5-7}{4} = -1/2
\mbox{\hspace{20pt}and\hspace{20pt}}
y=9\mbox{\hspace{10pt}becomes\hspace{10pt}}  z_y\frac{9-7}{4} = 1/2\]
\item Compute the products $z_x\,z_y$.
\item Compute the correlation coefficient.
\[r = \frac{1}{5}\,\left(\frac{3}{4} - \frac{1}{4} + 0 - \frac{3}{4} + \frac{9}{4}\right)
  = \frac{1}{5}\,\frac{8}{4} = \frac{2}{5}\]
\item Find the SD line.\vspace{-15pt}
\begin{align*}
\big(y-\mbox{mean}(y)\big) &= (\mbox{sign}\;\; r)\,\frac{\mbox{SD}_y}{\mbox{SD}_x}\,
   \big(x-\mbox{mean}(x)\big)\\
(y-7) & = \frac{4}{2}\,(x-4)\\
y-7   & = 2\,(x-4)\\
y     & = 2\,x -1
\end{align*}
\item Compute the $y$ values predicted by the SD line.   For example,
\[\mbox{if $x=1$, then $y=2\cdot 1 - 1=2 - 1 = 1$ \hspace{20pt} and\hspace{20pt} if $x=3$, then
  $y=2\cdot 3 - 1 = 6-1 = 5$.}\]
\item Compute the SD errors $(\mbox{error}=y - \mbox{predicted value of $y$})$.  For example,
\[\mbox{if $x=1$, then $\mbox{error}=5 - 1=4$\hspace{20pt} and\hspace{20pt} if $x=3$, then
  $\mbox{error}=9-5=4$.}\]
\item Compute the RMS size of the errors.
\[\mbox{RMS}_{\mbox{\scriptsize SD}} = \sqrt{(4^2 + 4^2 + (-8)^2)/5} = \sqrt{(16 + 16 + 64)/5}
  = \sqrt{96/5} \approx 4.38.\]
\end{itemize}
\vfill
\eject
{\ }

Here is how to get \texttt{R} to do all this:

\begin{tabular}{ll}
\lstinline!>source("http://www.adjoint-functors.net/SD.R")!\span\\
\lstinline!>source("http://www.adjoint-functors.net/SDline.R")!\span\\
\lstinline!> x <- c(1, 3, 4, 5, 7)!   & define the $x$ values\\
\lstinline!> y <- c(5, 9, 7, 1, 13)! & define the $y$ values\\[5pt]
\lstinline!> mean(x)!                 & We can actually skip from here $\dots$\\
\lstinline!> SD(x)!\\
\lstinline!> mean(y)!\\
\lstinline!> SD(y)!\\
\lstinline!> zx <-! \ \lstinline!zScore(x)! & \\
\lstinline!> zy <-! \ \lstinline!zScore(y)! & \\
\lstinline!> r <-! \ \lstinline!mean(zx*zy)! & \\
\lstinline!> SDslope <-! \ \lstinline! sign(r)*SD(y)/SD(x)! & \\
\lstinline!> yIntercept <-! \ \lstinline! mean(y) - SDslope*mean(x)! &  
   \hspace{30pt} $\dots$ to the next line\\[5pt]
\lstinline!> SDline(x, y)!               & Read the slope, $y$ intercept and $r$ from this\\
\lstinline!> predictedBySD <-! \ \lstinline!2*x - 1! & \\
\lstinline!> SDerrors <-! \  \lstinline!y - predictedBySD! & \\
\lstinline!> SDrms <-! \ \lstinline!sqrt(mean(SDerrors^2))! & \\[5pt]
\end{tabular}

\vfill
\eject

\end{document}

